\documentclass{article}
\usepackage{amsmath,amssymb,hyperref}
\title{Native bicolimit normalization and separated parallel}
\date{15 November 2026}
\begin{document}\maketitle
\paragraph{Scope.} A written transformation-level theorem. This is a weak, Cat-valued categorical lifting criterion retaining native modifications. It does not infer the criterion from Set-level connectedness, and it does not assert preservation by an arbitrary directed effectful parallel functor. Bicolimits, pseudonatural transformations and bicategorical Kan extensions are established machinery; no novelty claim is made.

\paragraph{Data.} Let $B,D$ be small bicategories and let $W,P:B\to\mathrm{Cat}$ and $Z,Q:D\to\mathrm{Cat}$ be pseudofunctors. Write $\operatorname{Ps}(B)(W,P)$ for the category of pseudonatural transformations and modifications. Assume specified bicolimit equivalences
\[
 \operatorname{bicolim}_B W\simeq1,
 \qquad\operatorname{bicolim}_D Z\simeq1.
\]
The assumption means, in particular, that the canonical constant comparison $X\to\operatorname{Ps}(D)(Z,\Delta X)$ is an equivalence pseudonaturally in the category $X$, and likewise for $W$. It is not an assertion only about connected components of objects.

\paragraph{Theorem.} Taking componentwise Cartesian products induces an equivalence of native categories
\[
 \operatorname{Ps}(B)(W,P)\times\operatorname{Ps}(D)(Z,Q)
 \simeq\operatorname{Ps}(B\times D)(W\boxtimes Z,P\boxtimes Q).
\]
It is natural in the indicated data with their specified pseudofunctorial comparisons, and retains all modifications rather than quotienting transformations to isomorphism classes.

\paragraph{Proof.} Split the target $P\boxtimes Q$ as $(P\pi_B)\times(Q\pi_D)$. Transformations and modifications into this product split into their two components. For the first component, curry a pseudonatural transformation $W_b\times Z\Rightarrow\Delta P_b$ at each $b$ to a functor
\[
 W_b\longrightarrow\operatorname{Ps}(D)(Z,\Delta P_b).
\]
The $D$-pseudonaturality constraints give the invertible constraints of each curried transformation; naturality in $W_b$ gives its modifications. The $B$-pseudonaturality constraints assemble these into a pseudonatural transformation from $W$ to $b\mapsto\operatorname{Ps}(D)(Z,\Delta P_b)$. Mixed-coordinate coherence is precisely the compatibility needed for this assembly. Conversely, evaluate the curried data and use a chosen coordinate factorization of arrows of $B\times D$. Pseudofunctor coherence gives inverse constructions up to canonical invertible modification. Thus
\[
 \operatorname{Ps}(B\times D)(W\boxtimes Z,P\pi_B)
 \simeq\operatorname{Ps}(B)(W,b\mapsto\operatorname{Ps}(D)(Z,\Delta P_b)).
\]
The bicolimit universal property identifies the inner category with
$\mathrm{Cat}(\operatorname{bicolim}Z,P_b)\simeq P_b$, pseudonaturally in $b$. This gives $\operatorname{Ps}(B)(W,P)$. Exchange the coordinates for the second component. Tracking the canonical constant comparisons shows that the resulting equivalence is the inverse of the displayed product operation. No strict Cartesian-closedness claim about the 2-category of strict 2-categories is used.

\paragraph{Right bi-Kan application.} For a pseudofunctor $i:B\to A$, use the pointwise right bi-Kan convention
\[
 (\operatorname{Ran}_iP)(a)\simeq
 \operatorname{Ps}(B)(A(a,i-),P).
\]
If every weight $A(a,i-)$ has bicolimit equivalent to $1$, and similarly for $j:D\to C$, the theorem yields
\[
 (\operatorname{Ran}_iP)\boxtimes(\operatorname{Ran}_jQ)
 \simeq\operatorname{Ran}_{i\times j}(P\boxtimes Q),
\]
with native modifications. The product hom weight is the external product of the two hom weights; the pointwise equivalences assemble by the same pseudonaturality. If $i$ has a left biadjoint $L$, then $A(a,i-)\simeq B(La,-)$ and bicategorical Yoneda verifies the normalization. This supplies a geometric sufficient condition, not a requirement that every comparison be representable.

\paragraph{Coherent finite iteration.} Products of normalized weights are again normalized: apply the curried universal property with constant target $X$ to obtain $\operatorname{Ps}(W\boxtimes Z,\Delta X)\simeq X$. Finite separated product comparisons are coherent because, on every component and modification, they are the same Cartesian tuple operation. Selected maps $i$ satisfying the weight condition are also closed under composition whenever the indicated pointwise Kan extensions exist: their right bi-Kan functors preserve all constant diagrams via the canonical comparison, this property composes, and the weighted universal property detects the normalization of the composite weights. These statements do not define a tensor on arbitrary coupled inputs or mixed new-arrow expressions.

\paragraph{Why Set connectedness is insufficient.} Let $B=BG$ for a nontrivial group and $W$ the constant terminal category. Set-valued constant cones see only connectedness. In contrast, pseudonatural transformations $W\Rightarrow\Delta X$ are functors $BG\to X$, with their natural transformations. Thus the bicolimit of $W$ is $BG$, not the terminal category. Ordinary strict Cat-enriched conical colimits may also discard the index coherence that this bicolimit retains. Consequently the earlier connected-lift/sifted-base result cannot simply be promoted to this native theorem; an appropriately stronger normalization must be proved.

\paragraph{Relation to the original relative factory.} The original construction is Cat-valued and retains directed cells and endpoint changes. The singleton-module action representation likewise produces pseudofunctors and pseudonatural old module maps, so this theorem has the correct native comparison level for selected reindexing/right-adjoint assembly. It supplies a structural criterion on the weights rather than a chosen physical model. It does not automatically make the original directed parallel operation preserve these limits: external Cartesian pairing of native categories is one input to that operation, not the effectful operation itself. Deriving the weight normalization and the required mixed compatibility from the original independent context/content data remains the next obligation.
\end{document}
